\documentclass[11pt]{article}

% ==== Paquetes ====
\usepackage[utf8]{inputenc}
\usepackage[T1]{fontenc}
\usepackage[spanish]{babel}
\usepackage{amsmath, amssymb, amsfonts}
\usepackage{graphicx}
\usepackage{caption}
\usepackage{subcaption}
\usepackage{float}
\usepackage{geometry}
\usepackage{hyperref}
\usepackage{minted}
\usepackage{mdframed}
\geometry{margin=2.3cm}

\usemintedstyle{friendly}
\setminted{fontsize=\scriptsize, breaklines=true}

% ==== Datos de la tarea ====
\title{IEE3864 – Fundamentos de instrumentación astronómica \\ \large Tarea 2}
\author{Nicolás Seguel}
\date{\today}

\begin{document}
\maketitle

\begin{enumerate}
    \item \textbf{given the definition, calculate the length of 1 pc in Km and light yr}
    
Un parsec (pc) es la distancia a la que la distancia entre la Tierra y el Sol (1 UA) subtendería un ángulo de 1 segundo de arco. Para calcular la longitud de 1 pc en kilómetros y años luz, podemos usar las siguientes fórmulas:

\[
1 \, \text{pc} = \frac{1 \, \text{AU}}{\tan(1 / 3600)}
\]

\[
1 \, \text{AU} = 149597870.7 \, \text{km}
\]

Calculando la longitud de 1 pc en kilómetros:

\[
1 \, \text{pc} = \frac{149597870.7 \, \text{km}}{\tan(1 / 3600)} \approx 3.0857 \times 10^{13} \, \text{km}
\]

Ahora en años luz se tiene que:

\[
1 \, \text{ly} = 300000000 \, \text{km/s} \times 365.25 \, \text{days} \times 24 \, \text{hours} \times 3600 \, \text{seconds}
\]

\[
1 \, \text{ly} \approx 9.461 \times 10^{12} \, \text{km}
\]

\[
1 \, \text{pc} = \frac{3.0857 \times 10^{13} \, \text{km}}{9.461 \times 10^{12} \, \text{km/ly}} \approx 3.2616 \, \text{ly}
\]

    \item \textbf{calculate the apparent size of a 30 Kpc object at a distance of 1000 Mpc}
    
Utilizando trigonometría, el tamaño aparente de un objeto se puede calcular como:

\[
\theta = \arctan\left(\frac{15 \, \text{kpc}}{1000 \, \text{Mpc}}\right) \times 2
\]

\[
\theta = 6,187 \, \text{arcsec}
\]

    \item \textbf{calculate the sun Irradiance at the Earth surface}
    
La irradiancia del Sol, también conocida como constante solar, es la cantidad de energía solar que llega a la superficie de la Tierra por unidad de área. Según Wikipedia, la constante solar se puede calcular con la siguiente fórmula:

\[
K = \sigma T^4 (\frac{r_s}{a_0})^2
\]

Donde:
\begin{itemize}
    \item $\sigma$ es la constante de Stefan-Boltzmann ($5.67 \times 10^{-8} \, \text{W/m}^2\text{K}^4$)
    \item $T$ es la temperatura del Sol (aproximadamente 5778 K)
    \item $r_s$ es el radio del Sol (aproximadamente $6.96 \times 10^8$ m)
    \item $a_0$ es la distancia media entre la Tierra y el Sol (aproximadamente $1.496 \times 10^{11}$ m)
\end{itemize}

Calculando la constante solar:

\[
K = 5.67 \times 10^{-8} \times (5778)^4 \times \left(\frac{6.96 \times 10^8}{1.496 \times 10^{11}}\right)^2
\]

\[K \approx 1361 \, \text{W/m}^2
\]

    \item \textbf{calculate the Irradiance of an object of $10^{12} L \odot$ at a distance of 1000 Mpc}
    
La irradiancia de un objeto se puede calcular utilizando la ley del inverso del cuadrado. La fórmula es:

\[
I = \frac{L}{4 \pi d^2}
\]

Donde:
\begin{itemize}
    \item $L$ es la luminosidad del objeto ($10^{12} L_\odot$)
    \item $d$ es la distancia al objeto (1000 Mpc)
    \item $L_\odot$ es la luminosidad del Sol (aproximadamente $3.828 \times 10^{26}$ W)
\end{itemize}

Calculando la luminosidad del objeto:

\[
I = \frac{10^{12} \times 3.828 \times 10^{26}}{4 \pi (1000 \times 3.086 \times 10^{19})^2}
\]

\[
I = 3.199 \times 10^{-14} \, \text{W/m}^2
\]

    \item \textbf{calculate the sun’s apparent magnitude in V}
    
La magnitud aparente del Sol en $m_v$ (magnitud visual) se puede calcular utilizando la fórmula:

\[
m_v = -2.5 \log_{10}\left(\frac{I}{I_0}\right)
\]

Donde:
\begin{itemize}
    \item $I$ es la irradiancia del Sol (1361 W/m²)
    \item $I_0$ es la irradiancia de referencia para la magnitud visual (aproximadamente $2.52 \times 10^{-8}$ W/m²)
\end{itemize}

Calculando la magnitud aparente del Sol:

\[
m_v = -2.5 \log_{10}\left(\frac{1361}{2.52 \times 10^{-8}}\right) \approx -26.74
\]

    \item \textbf{convert Jy in units of $Joule/sec/m^2/\mu m$}
    
Un Jansky (Jy) es una unidad de medida de la densidad de flujo espectral, que se define como:

\[
1 \, \text{Jy} = 10^{-26} \, \text{W/m}^2/\text{Hz}
\]

Para convertir Jy a $Joule/sec/m^2/\mu m$, necesitamos convertir la frecuencia (Hz) a longitud de onda ($\mu m$). La relación entre frecuencia y longitud de onda es:

\[
1 \, \text{Hz} = \frac{c}{\lambda}
\]

Reemplazando esta relación en la definición de Jy:

\[
1 \, \text{Jy} = 10^{-26} \, \text{W/m}^2 \times \frac{\lambda}{c}
\]

\[
1 \, \text{Jy} = 10^{-26} \times \frac{\lambda_{\mu m}}{3 \times 10^8} \, \text{Joule/sec/m}^2/\mu m
\]

\[1 \, \text{Jy} = 3.33 \times 10^{-35} \lambda_{\mu m} \, \text{Joule/sec/m}^2/\mu m
\]

\end{enumerate}

\end{document}